\documentclass[t, aspectratio=169]{beamer}
\usepackage{fontspec}
\setmainfont{Times New Roman}
\usepackage{amsmath}
\usepackage{amssymb}
\usepackage{booktabs}
\usepackage{graphicx}
\usepackage{xcolor}
\usepackage{array}
\usepackage{tikz}
\usetikzlibrary{shapes.callouts, positioning, arrows.meta, fit, backgrounds}

\definecolor{hustblue}{RGB}{0,60,113}
\definecolor{goodgreen}{RGB}{30,120,60}
\definecolor{warnorange}{RGB}{200,90,20}
\definecolor{badred}{RGB}{170,30,30}

\def\slideContentLeftMargin{0.5cm}
\def\slideContentRightMargin{0.5cm}
\def\hustHeaderHeight{3.0cm}
\def\slideContentBotMargin{2.8cm}

\usepackage[blue, 169]{beamerthemeHUST}
\usefonttheme{serif}
\renewcommand{\familydefault}{\rmdefault}

\setbeamercolor{frametitle}{fg=hustblue}
\setbeamerfont{frametitle}{size=\fontsize{15}{17}\selectfont, series=\bfseries}
\setbeamertemplate{frametitle}{
  \nointerlineskip
  \begin{tikzpicture}[remember picture, overlay]
    \node[anchor=north west, xshift=0.35cm, yshift=-0.95cm, text width=12.3cm, align=left]
    at (current page.north west)
    {\usebeamerfont{frametitle}\usebeamercolor[fg]{frametitle}\insertframetitle};
  \end{tikzpicture}
}
\setbeamertemplate{footline}{\vbox to 1.8cm{}}

\renewcommand{\titlePageTitleX}{-0.5cm}
\renewcommand{\titlePageTitleY}{-0.5cm}
\renewcommand{\hustTitleAlign}{\alignLeft}
\renewcommand{\hustSubtitleAlign}{\alignLeft}
\renewcommand{\hustAuthorAlign}{\alignLeft}
\renewcommand{\hustInstAlign}{\alignLeft}

\renewcommand{\hustsectionpage}[1]{%
  {%
    \setbeamertemplate{background}{\includegraphics[width=\paperwidth,height=\paperheight]{\hustbg{section}}}%
    \begin{frame}[plain,c]
      \centering \Huge \textcolor{white}{\textbf{#1}}
    \end{frame}
  }%
}

\graphicspath{{images/}}

\newcommand{\bignum}[4]{%
  \begin{minipage}[t]{#4}
    \centering
    {\fontsize{32}{36}\selectfont\bfseries\color{#3} #1}\\[2pt]
    {\small #2}
  \end{minipage}%
}
\newcommand{\bignumS}[4]{%
  \begin{minipage}[t]{#4}
    \centering
    {\fontsize{22}{26}\selectfont\bfseries\color{#3} #1}\\[1pt]
    {\scriptsize #2}
  \end{minipage}%
}

\title{\large Báo cáo tiến độ -- Tuần 1--2}
\subtitle{\normalsize Tích hợp dữ liệu khí tượng đa nguồn Đông Nam Á}
\author{Ngô Văn Quang \quad $\cdot$ \quad Nguyễn Cao Lương \quad $\cdot$ \quad Trương Tuấn Nghĩa}
\institute{SoICT, Đại học Bách khoa Hà Nội \\
IT5427 -- Tích hợp và xử lý dữ liệu lớn -- Nhóm 16\\[3pt]
{\small Giảng viên: TS.\ Vũ Tuyết Trinh \quad -- \quad 20/10/2026}}

\begin{document}

\husttitlepage

\begin{frame}[shrink=5]{Nội dung}
\begin{columns}[T]
\begin{column}{0.5\textwidth}
\begin{enumerate}
\item \hyperlink{sec1}{Tổng quan tiến độ 160 tác vụ}
\item \hyperlink{sec2}{Ingestion -- TV1}
\item \hyperlink{sec3}{Schema \& Warehousing -- TV2}
\end{enumerate}
\end{column}
\begin{column}{0.5\textwidth}
\begin{enumerate}
\setcounter{enumi}{3}
\item \hyperlink{sec4}{Spark \& Physics-Informed ML -- TV3}
\item \hyperlink{sec5}{Dashboard -- TV4}
\item \hyperlink{sec6}{Rủi ro \& kế hoạch Tuần 3}
\end{enumerate}
\end{column}
\end{columns}
\end{frame}

\resetPageCount
\showPageNumber
\useStyleHustFull
\useThinBar

%=============================================================
\section{Tổng quan tiến độ}
%=============================================================
\hypertarget{sec1}{}
\hustsectionpage{1. Tổng quan tiến độ\\160 tác vụ}

\begin{frame}{51/80 tác vụ Tuần 1--2 đã đạt đúng tiêu chí kiểm thử đã đặt ra từ đầu, không tự nhận suông}
\begin{center}
\bignum{51}{tác vụ \textbf{Xong}\\(đạt đủ Kiểm thử/DoD)}{goodgreen}{3.6cm}
\bignum{11}{tác vụ \textbf{Một phần}\\(đạt mục tiêu khác cách)}{warnorange}{3.6cm}
\bignum{18}{tác vụ \textbf{Chưa làm}\\(hành động của nhóm)}{badred}{3.6cm}
\end{center}
\vspace{0.3cm}
\centering\small
Mỗi trạng thái được chấm theo đúng cột ``Kiểm thử'' trong \texttt{docs/03-ke-hoach-160-tasks.md}
đã lập từ đề cương -- không phải tự đánh giá chủ quan. 18 tác vụ \textbf{Chưa làm} hầu hết là hành
động của con người (họp nhóm, nộp bài, diễn tập) ngoài khả năng tự động hoá.
\end{frame}

\begin{frame}{TV3+TV4 (kiêm nhiệm) hoàn thành tỉ lệ cao nhất nhờ ưu tiên đúng Spark/Dashboard trước PINN/Kafka}
\centering\small
\begin{tabular}{@{}lccc@{}}
\toprule
Vai trò & \textcolor{goodgreen}{\textbf{Xong}} & \textcolor{warnorange}{\textbf{Một phần}} & \textcolor{badred}{\textbf{Chưa làm}} \\
\midrule
TV1 -- Ingestion (Quang) & 6/10 & 3/10 & 1/10 \\
TV2 -- Schema (Lương) & 7/10 & 2/10 & 1/10 \\
TV3 -- Spark/PINN (Nghĩa) & 8/10 & 1/10 & 1/10 \\
TV4 -- Dashboard/PM (Nghĩa) & 4/10 & 1/10 & 5/10 \\
\bottomrule
\end{tabular}
\vspace{0.3cm}

\small TV4 thấp nhất vì 5 tác vụ còn lại là hoạt động nhóm tới hạn đúng hôm nay (họp, slide,
diễn tập, trình bày) -- không phải do thiếu code.
\end{frame}

%=============================================================
\section{Ingestion}
%=============================================================
\hypertarget{sec2}{}
\hustsectionpage{2. Ingestion\\(TV1)}

\begin{frame}{Pipeline đã nạp thật 1.076.473 bản ghi từ 5 nguồn, không còn là dữ liệu mẫu nhỏ}
\centering\small
\begin{tabular}{@{}lrl@{}}
\toprule
Nguồn & Số bản ghi thật & Phạm vi \\
\midrule
GHCN-Daily & 655.164 & 41 trạm (VN + ĐNA), toàn lịch sử tới 2025 \\
NOAA ISD/Global Hourly & 358.273 & 26 trạm, năm 2024 \\
Open-Meteo Historical & 52.704 & 6 điểm đại diện, năm 2024 \\
ERA5 (ARCO, nội suy Bilinear) & 8.136 & 113 trạm, 15--17/6/2021 \\
NASA POWER & 2.196 & 6 điểm đại diện, năm 2024 \\
\midrule
\textbf{Tổng} & \textbf{1.076.473} & \\
\bottomrule
\end{tabular}
\vspace{0.2cm}
\centering\small Nguồn số liệu: chạy thật \texttt{warehouse/build\_warehouse.py} 10/10/2026.
\end{frame}

\begin{frame}{GHCN và ISD dùng 2 hệ mã trạm độc lập cho cùng vị trí vật lý -- chưa thể so khớp qua ID}
\begin{center}
\begin{tikzpicture}[remember picture,
  box/.style={rectangle, rounded corners, draw=hustblue, thick, fill=hustblue!8, align=center, inner sep=6pt, font=\small},
  arr/.style={-{Latex[length=3mm]}, thick, badred}]
\node[box, minimum width=5cm] (g) at (-3.6,1) {GHCN\\\texttt{VMM00048820}};
\node[box, minimum width=5cm] (i) at (3.6,1) {ISD\\\texttt{489000-99999} (USAF+WBAN)};
\node[rectangle callout, callout absolute pointer={(0,0.6)}, fill=badred!10, text=badred,
  rounded corners, inner sep=6pt, align=center, font=\small]
  at (0,-1.3) {Cùng 1 trạm vật lý\\(Hà Nội) nhưng\\2 ID không khớp};
\draw[arr] (g) -- (0,-0.6); \draw[arr] (i) -- (0,-0.6);
\end{tikzpicture}
\end{center}
\vspace{0.1cm}
\centering\small
Cần Entity Resolution dựa trên \textbf{khoảng cách toạ độ} (không phải so khớp ID) --
đã ghi nhận làm việc Tuần 3 tại \texttt{docs/11} mục 1, không bỏ qua.
\end{frame}

%=============================================================
\section{Schema \& Warehousing}
%=============================================================
\hypertarget{sec3}{}
\hustsectionpage{3. Schema \&\\Warehousing (TV2)}

\begin{frame}{Kho Mediated Schema đạt 100\% referential integrity trên 119 vị trí, ETL không mất 1 bản ghi}
\begin{center}
\bignumS{119}{vị trí\\(113 trạm + 6 điểm)}{goodgreen}{3.5cm}
\bignumS{100\%}{referential\\integrity}{goodgreen}{3.5cm}
\bignumS{1.076.473=1.076.473}{bản ghi vào = ra\\qua Spark ETL}{goodgreen}{5cm}
\end{center}
\vspace{0.3cm}
\centering\small
\texttt{warehouse/location.py} + \texttt{warehouse/build\_warehouse.py}, partition theo
\texttt{year/source} qua Spark, 88 partition thật. Survivorship (ISD/GHCN $>$ ERA5 $>$ Open-Meteo
$>$ POWER) đã test với overlap thật GHCN vs ERA5 (15--17/6/2021): chọn đúng nguồn ưu tiên cao hơn.
\end{frame}

\begin{frame}{Partitioning theo year/source làm CHẬM truy vấn 2,2 lần ở quy mô dữ liệu hiện tại}
\begin{center}
\begin{tikzpicture}[remember picture,
  bar/.style={fill=hustblue!70}, barslow/.style={fill=badred!70}]
\draw[-{Latex[length=2.5mm]}, thick] (0,0) -- (0,3.2) node[above] {\scriptsize giây};
\draw[-{Latex[length=2.5mm]}, thick] (0,0) -- (7,0);
\draw[bar] (1,0) rectangle (2.5,1.73*1.5) node[midway, above, font=\small] {};
\node[font=\small] at (1.75,1.73*1.5+0.25) {0,173s};
\node[font=\scriptsize] at (1.75,-0.35) {Không partition};
\draw[barslow] (4,0) rectangle (5.5,3.81*1.5);
\node[font=\small] at (4.75,3.81*1.5+0.25) {\textbf{0,381s}};
\node[font=\scriptsize] at (4.75,-0.35) {Có partition (88 thư mục)};
\node[rectangle callout, callout absolute pointer={(4.75,2.9)}, fill=badred!10, text=badred,
  rounded corners, inner sep=5pt, align=center, font=\scriptsize]
  at (6.3,3.4) {Chậm hơn\\$\approx$2,2 lần};
\end{tikzpicture}
\end{center}
\vspace{0.1cm}
\centering\small
Nguyên nhân: ở quy mô $\sim$1 triệu dòng, Spark quét toàn bộ cực nhanh; overhead liệt kê 88 thư
mục nhỏ (small-file problem) vượt lợi ích partition pruning. Giữ nguyên partition key, benchmark
lại khi ERA5 full 5 năm ($\approx$35GB) -- \texttt{docs/10} mục 2.
\end{frame}

%=============================================================
\section{Spark \& PINN}
%=============================================================
\hypertarget{sec4}{}
\hustsectionpage{4. Spark \&\\Physics-Informed ML (TV3)}

\begin{frame}{Ma trận tương quan Spark xác nhận gió-độ ẩm là cặp biến liên kết mạnh nhất (+0,57)}
\centering\footnotesize
\begin{tabular}{@{}lrrrrr@{}}
\toprule
 & temp\_c & rh\_pct & precip\_mm & wind\_speed & pressure\_msl \\
\midrule
temp\_c & \textbf{1,00} & 0,11 & -0,07 & 0,16 & 0,06 \\
rh\_pct & 0,11 & \textbf{1,00} & -0,16 & \textcolor{goodgreen}{\textbf{0,57}} & 0,43 \\
precip\_mm & -0,07 & -0,16 & \textbf{1,00} & -0,10 & -0,07 \\
wind\_speed & 0,16 & \textcolor{goodgreen}{\textbf{0,57}} & -0,10 & \textbf{1,00} & 0,51 \\
pressure\_msl & 0,06 & 0,43 & -0,07 & 0,51 & \textbf{1,00} \\
\bottomrule
\end{tabular}
\vspace{0.25cm}
\centering\small
Tính thật trên 1.076.473 dòng (\texttt{VectorAssembler} + \texttt{Correlation.corr}). Đường chéo
=1,0, ma trận đối xứng -- đúng kiểm thử. Gió mạnh đi kèm hệ thống áp thấp/mưa -- hợp lý khí tượng.
\end{frame}

\begin{frame}{Ba kỹ thuật Physics-Informed Neural Network đều hội tụ đúng nghiệm vật lý kỳ vọng}
\centering\small
\begin{tabular}{@{}p{3.8cm}p{4cm}p{4cm}@{}}
\toprule
Kỹ thuật & Kiểm chứng & Kết quả thật \\
\midrule
Phương trình nhiệt 1D & So nghiệm giải tích $e^{-\alpha\pi^2 t}\sin(\pi x)$ & Hội tụ đúng sau 100 epoch \\
Clausius-Clapeyron (RH từ điểm sương) & Physics loss trên 5.000 bản ghi ISD thật & Loss 1068 $\to$ 2,9, MAE 1,07 điểm\% \\
Navier-Stokes 2D rút gọn & Gradient qua \texttt{autograd} không lỗi & Phát hiện đúng 1 case \texttt{grad=None} hợp lệ (bias lớp cuối) \\
\bottomrule
\end{tabular}
\vspace{0.2cm}
\centering\small \texttt{model/train\_demo.py}, \texttt{model/pinn\_rh\_from\_dewpoint.py},
\texttt{model/navier\_stokes\_autograd\_check.py} -- 3/3 test subprocess pass.
\end{frame}

\begin{frame}{Cụm Spark 3 node và Kafka broker đã chạy THẬT hôm nay, không còn là kế hoạch trên giấy}
\begin{center}
\bignumS{Kafka}{broker thật, producer\\publish 70 message JSON}{goodgreen}{4.5cm}
\bignumS{3 node}{Spark cluster\\(1 master + 2 worker)}{goodgreen}{4.5cm}
\end{center}
\vspace{0.2cm}
\centering\small
\texttt{docker-compose.kafka.yml} + \texttt{kafka-console-consumer} đọc lại đúng JSON từ topic
\texttt{weatherobs.ghcn}. Phát hiện \& sửa 2 lỗi thật ngay trong lúc triển khai: (1) tag
\texttt{bitnami/kafka:3.7}/\texttt{spark:3.5} bị Bitnami rút khỏi Docker Hub -- đổi sang
\texttt{bitnamilegacy}; (2) producer xuất \texttt{NaN} (không hợp lệ JSON chuẩn) cho giá trị
thiếu -- đã sửa thành \texttt{null}.
\end{frame}

%=============================================================
\section{Dashboard}
%=============================================================
\hypertarget{sec5}{}
\hustsectionpage{5. Dashboard\\(TV4)}

\begin{frame}{Dashboard Streamlit hiển thị dữ liệu thật qua 3 thành phần, không còn dữ liệu giả}
\centering\small
\begin{tabular}{@{}p{3.6cm}p{8.4cm}@{}}
\toprule
Thành phần & Trạng thái thật \\
\midrule
Chuỗi thời gian theo trạm & Dropdown 73 trạm/điểm thật, đổi trạm cập nhật đúng dữ liệu \\
Bảng tổng hợp theo nguồn & Tổng hợp trực tiếp từ warehouse Parquet (pandas, không mock) \\
Bản đồ + time slider & \texttt{st.map()}, join với dimension \texttt{Location} (119 vị trí), mặc định 15/6/2024 (66 vị trí hợp lệ) \\
\bottomrule
\end{tabular}
\vspace{0.2cm}
\centering\small
\texttt{dashboard/app.py} -- đã chạy sạch \texttt{streamlit run}, verify bằng mắt không lỗi runtime.
\end{frame}

%=============================================================
\section{Rủi ro \& kế hoạch}
%=============================================================
\hypertarget{sec6}{}
\hustsectionpage{6. Rủi ro \&\\kế hoạch Tuần 3}

\begin{frame}{Bốn rủi ro mới phát hiện khi chạy pipeline thật đều đã có phương án dự phòng cụ thể}
\centering\scriptsize
\begin{tabular}{@{}p{4.3cm}p{1.6cm}p{6.1cm}@{}}
\toprule
Rủi ro mới & Mức ảnh hưởng & Phương án dự phòng \\
\midrule
GHCN/ISD khác hệ mã trạm, chưa so khớp được qua ID & Trung bình-Cao & Entity Resolution theo khoảng cách toạ độ, Tuần 3 \\
Partition hiện tại chậm \& tốn dung lượng hơn ở quy mô 1M dòng & Thấp (hiện tại) & Giữ partition key, benchmark lại khi ERA5 full 5 năm ($\approx$35GB) \\
ERA5 mới có 3 ngày, chưa đủ 5 năm như đề cương & Trung bình & Quyết định phạm vi tải phụ thuộc thời gian/băng thông -- cần nhóm chốt \\
Khối lượng TV3+TV4 kiêm nhiệm vẫn nặng nhất & Cao & Đã ưu tiên đúng: Spark batch + Dashboard trước, PINN/Kafka sau \\
\bottomrule
\end{tabular}
\end{frame}

\begin{frame}{Tuần 3 tập trung đúng 3 việc: Entity Resolution, ERA5 quy mô lớn, và Spark Streaming thật}
\begin{center}
\bignumS{Entity\\Resolution}{GHCN $\leftrightarrow$ ISD\\theo khoảng cách toạ độ}{hustblue}{4cm}
\bignumS{ERA5\\quy mô lớn}{quyết định phạm vi\\tải 5 năm ($\approx$35GB)}{hustblue}{4cm}
\bignumS{Spark\\Streaming}{windowing thật trên\\Kafka broker đã chạy}{hustblue}{4cm}
\end{center}
\vspace{0.3cm}
\centering\small
51/80 tác vụ Tuần 1--2 đã Xong đúng tiêu chí kiểm thử; hạ tầng Docker (Spark cluster + Kafka)
đã chạy thật ngay trong ngày báo cáo -- sẵn sàng cho Tuần 3 không còn vướng hạ tầng.
\end{frame}

\hustthankyou

\end{document}
